% --------------------------------------------------
% Proposta de Trabalho (convertida de Markdown)
% --------------------------------------------------
\section{Proposta de Trabalho}

\subsection{Dataset}

\subsubsection{ICNF}
O principal conjunto de dados consiste numa versão filtrada dos registos do repositório \texttt{icnf\_fire\_data}, que agrega ocorrências do serviço \texttt{fogos.icnf.pt}. Para este estudo serão considerados apenas os incêndios entre 2020 e 2025, após remoção de registos irrelevantes, nomeadamente:
\begin{itemize}
	\item falsos alarmes;
	\item registos com coordenadas em falta;
	\item incêndios com área ardida inferior a 5 hectares;
	\item queimadas controladas;
	\item ocorrências não classificadas como incêndios;
	\item registos sem data/hora de início ou fim.
\end{itemize}

O ano de 2025 inclui dados até 06/10/2025.

\subsubsection{Meteorologia}
Foi desenvolvido um script Python que consulta a API do \texttt{open-meteo}. Para cada ano são gerados dois ficheiros:
\begin{itemize}
	\item \texttt{meteorology\_[ano].csv} — registos diários de meteorologia por \texttt{incident\_id};
	\item \texttt{summary\_[ano].csv} — agregação por \texttt{incident\_id} (médias, contagens e coordenadas), usada para validação.
\end{itemize}

\subsubsection{RSS Google News}
Foram recolhidos artigos noticiosos usando filtros por idioma (português), termos-chave e período (2020–2025). Extraíram-se metadados (título, data, fonte, URL). Ainda não foi definida a utilização final destes dados, mas visam complementar informação sobre meios e cobertura mediática.

\subsubsection{DECIR}
Foram identificados PDFs da Diretiva Operacional Nacional nº2 (2020–2025) com informação sobre meios aéreos e terrestres mobilizados. A extração automática ainda não foi concluída devido à formatação irregular dos PDFs.

\subsubsection{MODIS — MOD13Q1 (NDVI)}
Será usado o produto MOD13Q1 (NDVI 16-day, 250 m) através do Google Earth Engine. Para cada perímetro de incêndio serão calculadas estatísticas zonais pré e pós-incêndio (média, contagem de píxeis, etc.). O período foco é 2020–2024.

\subsubsection{CORINE Land Cover}
Usado para atribuir classes de uso do solo a perímetros de incêndio e para construir áreas de controlo. O produto base será o CLC 2018.

\subsubsection{INE}
Indicadores socioeconómicos a nível municipal/freguesia (demografia, emprego, sector agrícola, turismo) a integrar para avaliar impactos locais e vulnerabilidades.

\subsection{Objetivos}
\begin{enumerate}
	\item Identificar áreas com maior risco de incêndio;
	\item Determinar épocas do ano com maior ocorrência de incêndios;
	\item Identificar anos com maior número e gravidade de incêndios;
	\item Analisar o impacto das temperaturas no número, gravidade e duração dos incêndios;
	\item Avaliar a influência do estado da vegetação (NDVI) na probabilidade e severidade dos incêndios;
	\item Avaliar o efeito do tipo de uso do solo na ocorrência e extensão dos incêndios;
	\item Examinar como meios terrestres e aéreos afetam duração e gravidade;
	\item Avaliar impactos nos níveis de freguesia.
\end{enumerate}

\subsection{Questões Analíticas}
Apresentam-se pontos de análise associados aos objetivos (exemplos resumidos):
\begin{itemize}
	\item Incêndios por zona/distrito/concelho; área ardida total; percentagem de área ardida.
	\item Sazonalidade: incêndios por mês/estação, área média por mês.
	\item Tendências anuais e classificação por severidade.
	\item Relação entre faixas de temperatura (média, máxima) e métricas de incêndio (contagem, área, duração).
	\item Correlações entre NDVI pré-incêndio e probabilidade de ocorrência; comparações com áreas de controlo.
	\item Distribuição de incêndios por classe CORINE; área média por classe.
	\item Métricas de eficácia dos meios (tempo de resposta, meios mobilizados, taxa de expansão antes/depois da intervenção).
	\item Indicadores socioeconómicos pré e pós-incêndio a nível de freguesia (população, emprego, turismo, valores imobiliários).
\end{itemize}

\subsection{Identificação dos Processos de Negócio}
\begin{itemize}
	\item \textbf{Ocorrência de Incêndio}: registo do evento (localização, duração, área, causa).
	\item \textbf{Meteorologia Diária}: registo de condições meteorológicas por dia/ local.
	\item \textbf{Operações de Combate}: registo de meios mobilizados, tempos e eficácia.
	\item \textbf{Monitorização de Vegetação}: medições periódicas de índices (NDVI) e classificação.
\end{itemize}
\subsection{Metodologia}
Apresenta-se, de forma concisa, a abordagem metodológica seguida para modelação dimensional e desenho do processo ETL. Adotámos a abordagem orientada aos processos de negócio (Kimball), identificando o grão, as principais dimensões e as medidas para cada processo. Abaixo descrevem-se os quatro processos identificados, cada um com o respetivo grão e medidas principais.

\subsubsection{Ocorrência de Incêndio}
Grão: cada linha representa um incêndio (evento) identificado pelo \texttt{incident\_id}.

Principais dimensões: tipo de incêndio, entidade regista (ICNF), data/hora de início e fim, localização (coordenadas, freguesia, concelho, distrito).

Medidas principais: área ardida (ha), duração (horas), número de focos iniciais. Medidas derivadas a considerar: severidade (índice composto) e taxa de expansão (ha/h).

\subsubsection{Meteorologia Diária}
Grão: registo diário de condições meteorológicas associado a uma localização ou a um perímetro de incêndio.

Medidas principais: temperatura média e máxima (ºC), humidade relativa (%), velocidade do vento (km/h), precipitação (mm). Estes valores servem para calcular indicadores derivados (ex.: número de dias com temperatura acima de um limiar, soma de precipitação no período do incêndio).

\subsubsection{Operações de Combate}
Grão: cada registo corresponde a uma operação ou intervenção num incêndio específico.

Medidas principais: número de operacionais mobilizados, veículos terrestres, meios aéreos, tempo de resposta inicial (min), tempo total de operação/contensão e eficácia estimada da contenção. Serão também registadas datas/horas de chegada dos meios para análises temporais.

\subsubsection{Monitorização de Vegetação}
Grão: medições de índice de vegetação (por ex., NDVI) agregadas por perímetro e por janela temporal (pré/post-incêndio).

Medidas principais: NDVI médio pré-incêndio, NDVI médio pós-incêndio, contagem de píxeis validos, variação percentual. Estas medidas serão usadas em análises comparativas e como variáveis explicativas em modelos estatísticos.

\subsubsection{Notas sobre modelação e ETL}
O desenho do Data Warehouse segue princípios dimensionais: tabelas de factos centradas nos processos acima e dimensões partilhadas (tempo, local, tipo de uso do solo — CORINE, e entidades). O processo ETL inclui etapas de validação (remoção de registos com coordenadas em falta, filtro por área mínima, padronização de datas), enriquecimento (junção com meteorologia e NDVI) e agregação (fichas resumo por incident\_id).

Esta organização pretende facilitar a leitura num formato de artigo, concentrando a discussão metodológica numa subseção sucinta e orientada para as decisões de modelação.
